% ===================== FRONTMATTER =====================
\pagestyle{corporativo}

\begin{frontmatter}

	\title{\projecttitle}

	\author[inst1]{J.Arboleda\corref{cor1}}
	\ead{proyectos2@dmlsas.com}

	\author[inst1]{F. Navia}
	\ead{fernando.navia@dmlsas.com}

	\author[inst2]{H. Rosero}
	\ead{herminsul.rosero@dmlsas.com}

	\cortext[cor1]{Autor correspondiente}
	\address[inst1]{DML Ingenieros consultores, Cali, Colombia, 760025}
	\address[inst2]{Especialidad Mecánica y procesos, Cali, Colombia}

	\begin{abstract}
		\thispagestyle{corporativo}
		Este informe presenta el cálculo termodinámico detallado para una caldera acuotubular bagacera con capacidad de producción de 100~t/h de vapor sobrecalentado a \SI{106}{barg} y \SI{545}{\celsius}. Se emplea bagazo de caña de azúcar como combustible, con una humedad del 48\% y un contenido de cenizas del 10\% en base húmeda (AR), valor que contempla la materia extraña mineral incorporada durante la cosecha mecanizada sin lavado previo, constituyendo un escenario conservador de diseño. Las propiedades termodinámicas del agua y vapor se calcularon con las formulaciones IAPWS-97. Se determinó un consumo de bagazo de 37.67~t/h, correspondiente a un ratio de \textbf{2.655 toneladas de vapor por tonelada de bagazo}. Adicionalmente, se incluye una guía paso a paso para la simulación del proceso en Aspen Plus.
	\end{abstract}

	\begin{keyword}
		caldera bagacera \sep bagazo de caña \sep balance térmico \sep vapor sobrecalentado \sep cogeneración \sep Aspen Plus
	\end{keyword}

\end{frontmatter}

\pagestyle{corporativo}
\thispagestyle{corporativo}
